\subsection{A sharper first-unit bound in the special Cohen weight}

This argument improves the discriminant bound for the automatic family
$b=(p+3)/2$.  It supplies the class-number unit and the first twisted
Bernoulli unit.  It does not supply the second twisted unit
$\beta_3(\chi_D)$ required by the same-discriminant detector.

\begin{theorem}[KT-COHEN-WEIGHT-TWO-BOUND]
Let $p\geq31$ be prime with $p\equiv3\pmod4$, and set
$b=(p+3)/2$.  There exists an odd negative fundamental discriminant
$D=-d$ such that
\[
 3\leq d\leq p-4,\qquad
 p\nmid h(D),\qquad \beta_b(\chi_D)\in\mathbf Z_p^\times.
 \tag{CW.1}
\]
In particular $p\nmid D$.  By the already checked character transport,
$B_{1,\chi_D}$ and $B_{1,\chi_D\omega^{b-1}}$ are units, and
$\beta_n(\chi_D)$ is a unit for every odd $n\geq3$ congruent to $b$
modulo $p-1$.
\end{theorem}

\begin{proof}
Write $q=\exp(2\pi i\tau)$ and use the Cohen normalization
\[
 H_b(\tau)=\sum_{n\geq0}H(b,n)q^n
       \in M_{b+1/2}^{+}(\Gamma_0(4)).
 \tag{CW.2}
\]
The index-one Jacobi form corresponding to it is
$\zeta(1-2b)E_{b+1,1}$, as in KT-JACOBI-TWO-UNIT-FIELD.
Its full coefficient formula and the integrality proof there show
that every $H(b,n)$ is $p$-integral.  This includes indices divisible
by $p$: here $3\leq b\leq p-2$ and $b\ne(p-1)/2$.
The constant and support are
\[
 H(b,0)=-\frac{B_{p+3}}{p+3}\equiv\frac1{120}\pmod p,
 \qquad H(b,n)=0\quad(n\equiv1,2\pmod4).
 \tag{CW.3}
\]
The constant congruence is the same Kummer congruence used in
KT-COHEN-SPECIAL-HEAT-FILTRATION; in particular it is a unit.

Put
\[
 \theta(\tau)=\sum_{a\in\mathbf Z}q^{a^2},\qquad
 \Theta=\theta^4,
 \qquad G=2E_2(2\tau)-E_2(\tau).
 \tag{CW.4}
\]
Here $G\in M_2(\Gamma_0(2))\subset M_2(\Gamma_0(4))$ and
$\Theta\in M_2(\Gamma_0(4))$, both with trivial character and
integral Fourier coefficients.  The transformation law of $E_2$
shows the asserted modularity of $G$: its two nonmodular terms
cancel.  Their first coefficients are
\[
 \Theta=1+8q+24q^2+O(q^3),\qquad
 G=1+24q+24q^2+O(q^3).
 \tag{CW.5}
\]
The theta product
$\theta=\prod_{n\geq1}(1-q^{2n})(1+q^{2n-1})^2$
has no zeros in the upper half-plane.  The three cusps of
$\Gamma_0(4)$ are $\infty,0,1/2$.  At $\infty$ the theta
constant is $1$.  At $0$, Poisson summation gives
\[
 (\Theta|_2S)(\tau)=-\frac14\theta(\tau/4)^4,
 \qquad S=\begin{pmatrix}0&-1\\1&0\end{pmatrix},
 \tag{CW.6}
\]
so its constant is nonzero.  At $1/2$, with
$A=\left(\begin{smallmatrix}1&0\\2&1\end{smallmatrix}\right)$,
the same Poisson formula, with consistent holomorphic square roots,
is
\[
 (2\tau+1)^{-1/2}\theta(A\tau)
      =\sum_{a\text{ odd}}q^{a^2/4}
      =2q^{1/4}(1+O(q^2)).
 \tag{CW.7}
\]
Consequently $\Theta|_2A=16q+O(q^3)$; its sole zero is
the simple zero at this cusp.  The cusp has width $1$.
Also $G|_2A=G$, because $A\in\Gamma_0(2)$, so its constant
at that cusp is exactly $1$.

Now form the meromorphic modular form
\[
                       F=H_b/\theta^p.
 \tag{CW.8}
\]
Its weight is $b+1/2-p/2=2$.  More precisely the theta
automorphy factor of $H_b$ has exponent $2b+1=p+4$;
division by $\theta^p$ leaves the exponent $4$, which is
exactly the trivial-character weight-two factor on $\Gamma_0(4)$.
Thus there is no residual character in (CW.8).
This form is holomorphic away from the possible pole at $1/2$.
Its series at $\infty$ is $p$-integral, since $\theta$ has
constant $1$ and integral coefficients.  Frobenius gives
$\theta^p\equiv\theta(q^p)\pmod p$.  Therefore
\[
 [q^0]F\equiv1/120,
 \qquad [q]F\equiv[q^2]F\equiv0\pmod p.
 \tag{CW.9}
\]

We record explicitly which coefficients control its pole.
The two-component theta decomposition is
\[
 H_b(\tau)=h_0(4\tau)+h_1(4\tau),\qquad
 h_1(\tau)=\sum_{n\equiv3\ (4)}H(b,n)q^{n/4}.
 \tag{CW.10}
\]
For $v=\tau+1/2$ one has $4A\tau=2-1/v$.
The dual index-one theta representation has
$T^2=\operatorname{diag}(1,-1)$ and its $S$ matrix is an
eighth root of unity times
$2^{-1/2}\left(\begin{smallmatrix}1&1\\1&-1\end{smallmatrix}\right)$.
Multiplying its row $(1,-1)$ by that $S$ matrix leaves a
unit multiple of $(0,\sqrt2)$.  Hence, for an eighth root
of unity $u$ whose precise value is immaterial here,
\[
 (2\tau+1)^{-b-1/2}H_b(A\tau)
               =u\,2^{-b}h_1(\tau+1/2).
 \tag{CW.11}
\]
This derivation fixes both the component and the power of $2$;
all scalar factors are $p$-adic units.
Dividing (CW.11) by the $p$th power of (CW.7) shows that
$F|_2A$ has a $p$-integral Laurent expansion in integral
powers of $q$.  Its possible negative powers arise only
from $H(b,n)$ with $0<n<p$ and $n\equiv3\pmod4$.
Indeed its numerator terms have exponent $n/4$, whereas
its denominator is $2^p q^{p/4}$ times a power-series unit.
Since $n\equiv p\pmod4$, the resulting exponents are integral.
Its pole order is at most $(p-3)/4$.

Suppose, for contradiction, that all these odd-index coefficients
are zero modulo $p$.  Then every coefficient in the principal
part of $F|_2A$ is divisible by $p$.
The following explicit removal of that principal part justifies
the reduction step without an assumed base-change theorem for
half-integral weight forms or a special-fiber divisor assertion.
Work in a number field containing $u$ and the eighth roots of unity,
localized at a prime above $p$.  For every $j\geq1$ put
\[
                  W_j=G^{j+1}/\Theta^j.
 \tag{CW.12}
\]
These are weight-two meromorphic forms with trivial character.
They are holomorphic away from $1/2$, are integral at $\infty$,
and satisfy
\[
 W_j|_2A=16^{-j}q^{-j}+\text{higher powers of }q.
 \tag{CW.13}
\]
All Laurent coefficients here are integral at $p$, by
(CW.5) and (CW.7).  Start with the largest pole order of $F$.
Its coefficient, multiplied by $16^j$, is an element divisible
by $p$.  Subtract this multiple of $W_j$.  The leading pole
disappears; all remaining negative coefficients stay divisible
by $p$, because $W_j|_2A$ is integral.  Repeating finitely many
times yields a holomorphic characteristic-zero form
\[
                  F_0\in M_2(\Gamma_0(4))
 \tag{CW.14}
\]
whose expansion at $\infty$ is integral and congruent to that
of $F$ modulo $p$.  Every subtracted multiple is divisible
by $p$ and each $W_j$ is integral there.  Holomorphy at the
other cusps and in the upper half-plane was never changed.

The ordinary Sturm bound at weight $2$, level $4$ is
$2[\mathrm{SL}_2(\mathbf Z):\Gamma_0(4)]/12=1$.
Thus an integral form modulo $p$ is determined by its constant
and its $q$ coefficient.  The constant/$q$ coefficient matrix
of $\Theta,G$ has determinant $16$, a unit.  There are integral
scalars $a,c$ whose linear combination $a\Theta+cG$ has the
same first two coefficients as $F_0$.  By the Sturm theorem
they agree modulo $p$ at every coefficient.  Its $q$ and
$q^2$ coefficients would consequently satisfy
\[
 8a+24c=0,\qquad24a+24c=0\pmod p.
 \tag{CW.15}
\]
The determinant is $-384$, a unit for $p\geq31$, so
$a=c=0$ modulo $p$.  This contradicts the constant
$a+c\equiv1/120$ in (CW.9).  The contradiction proves that
some $H(b,n)$ with $0<n<p$, $n\equiv3\pmod4$, is a unit.

Choose the smallest such $n$.  Write $-n=D f^2$ with
$D<0$ fundamental.  It is odd, and $|D|\leq n<p$.
The full Cohen coefficient formula is
\[
 H(b,n)=-\beta_b(\chi_D)
       \sum_{v\mid f}\mu(v)\chi_D(v)v^{b-1}
                                  \sigma_{2b-1}(f/v).
 \tag{CW.16}
\]
The sum is an integer and the Bernoulli factor is $p$-integral.
A unit left side therefore forces $\beta_b(\chi_D)$ to be
a unit.  If $f>1$, the smaller fundamental coefficient
$H(b,|D|)=-\beta_b(\chi_D)$ would be a unit, contradicting
minimality.  Thus $n=d=|D|$ is fundamental.
Since $p$ and $d$ are both $3$ modulo $4$ and $d<p$,
one has $d\leq p-4$.
The carried elementary class-number estimate
KT-SMALL-DISCRIMINANT-CLASS-BOUND applies since $d<4p$;
it gives $h(D)<p$.  This proves (CW.1).
The stated character and all-weight consequences are exactly
the previously checked Kummer transport at conductor prime to $p$.
\end{proof}

\paragraph{Exact inputs and scope.}
The Cohen/Jacobi coefficient formula and its $p$-integrality,
including the conductor-divisible-by-$p$ case, are the already
whole-reviewed proof of KT-JACOBI-TWO-UNIT-FIELD, using
Eichler--Zagier's index-one theta decomposition and Theorem~2.1
(1983 primary preprint, printed pp.~19a--20 for the full
coefficient formula).  The two-component $S,T$ calculation is
displayed in (CW.10)--(CW.11); the theta Poisson normalization is
the same as the whole-reviewed KT-KOHNEN-THETA-NONCONGRUENCE.
The ordinary elliptic Sturm theorem is used with the exact
coefficient-field generality stated by Choi--Choie--Kikuta
(2013), printed p.~129.  These are imported modular-form inputs;
no new foundational proof of the Sturm theorem is claimed.
The finite principal-part removal proves the needed integral
reduction directly, so the argument does not require an
unverified identification of a special-fiber half-integral
weight space with reductions of characteristic-zero forms.

The new discriminant bound is linear in $p$.  It does not put
$d$ in the sufficient window $d^5\leq1600p^2$ of
KT-BETA3-SHORT-WINDOW.  No unit assertion for
$\beta_3(\chi_D)$, no full three-unit auxiliary field, and
no new vanishing of an original rational class-group component
is inferred from (CW.1).
